\chapter{Angiogram}

\section*{Definition and Overview}
An angiogram is an imaging test that uses X-rays and an injected contrast dye to examine the inside of blood vessels. A thin, flexible tube called a catheter is guided into an artery or vein, usually through a small puncture in the groin or wrist, and dye is injected while a sequence of X-ray images is recorded. The dye makes the vessels visible, revealing narrowing, blockages, aneurysms, abnormal connections, or sites of bleeding.

Coronary angiography is the most common form and examines the arteries supplying the heart, but the same technique is used to study the brain, neck, lungs, kidneys, and limbs. Angiography is both a diagnostic tool and a platform for treatment, because instruments such as balloons, stents, and coils can be delivered through the same catheter to treat the abnormality that has been found. Modern practice increasingly uses non-invasive alternatives such as CT and MRI angiography where appropriate, but catheter angiography remains the reference standard for many vascular questions.

\section*{Epidemiology}
Angiography is among the most common hospital procedures in the world, and coronary angiography is the most frequently performed type, with millions of procedures undertaken each year. It is used mainly in people with known or suspected coronary artery disease, and so its use increases with age and with the prevalence of vascular risk factors. Peripheral angiography of the legs, renal angiography, and cerebral angiography are performed in smaller numbers but are routine in specialist vascular, renal, and stroke services.

The pattern of use has evolved: purely diagnostic catheter studies have partly been replaced by CT and MRI angiography for initial assessment, while catheter angiography is increasingly reserved for cases where treatment is anticipated in the same session. As a result, a growing proportion of procedures combine diagnosis with intervention.

\section*{Aetiology and Pathophysiology}
An angiogram is a test rather than a disease, and the reasons for performing one relate to disorders of the blood vessels:

\begin{itemize}
    \item \textbf{Atherosclerosis:} fatty plaques narrowing or occluding the arteries of the heart, neck, kidneys, or legs
    \item \textbf{Aneurysms:} abnormal ballooning of an arterial wall that requires measurement, monitoring, or treatment
    \item \textbf{Thromboembolism:} blood clots or emboli obstructing vessels in the lungs, brain, or limbs
    \item \textbf{Bleeding:} localisation of the source of internal bleeding after injury, surgery, or in conditions such as gastrointestinal haemorrhage
    \item \textbf{Congenital anomalies:} abnormal vascular connections or malformations present from birth
    \item \textbf{Vasculitis:} inflammatory disease of the vessel walls
    \item \textbf{Pre-procedural planning:} precise measurement of blockages or anatomy before angioplasty, stenting, or surgery
\end{itemize}

The images obtained rely on the principle that blood, and the dye dissolved in it, absorbs X-rays less than the surrounding tissues; after injection, the dye column outlines the lumen of the vessel, and any irregularity, narrowing, or occlusion is displayed in real time.

\section*{Clinical Features}
An angiogram itself causes few symptoms, and what the patient feels reflects both the procedure and the condition being investigated:

\begin{itemize}
    \item Numbness and mild discomfort at the puncture site from the local anaesthetic
    \item A warm, flushing sensation spreading through the body as the contrast dye is injected
    \item A sense of pressure or brief awareness of the catheter moving, although the inside of the vessels has little sensation
    \item Mild bruising or a small lump (haematoma) at the puncture site afterwards
    \item The procedure typically takes 30--60 minutes and is performed under local anaesthetic, usually with mild sedation
    \item The findings may explain presenting symptoms such as chest pain, breathlessness, leg cramping, or stroke-like events
\end{itemize}

Most people are fully awake during the procedure and can be discharged home the same day or the next morning, depending on the access site and the results.

\section*{Differential Diagnosis}
An angiogram is performed when the history, examination, or other tests suggest vessel disease, and it helps to distinguish between:

\begin{itemize}
    \item Coronary artery narrowing versus other causes of chest pain
    \item Occlusive peripheral arterial disease versus nerve, joint, or venous causes of leg pain
    \item A true aneurysm versus a tortuous but normal vessel
    \item Pulmonary embolism versus other causes of breathlessness and chest pain
    \item An actively bleeding source versus a healed or non-bleeding lesion
    \item Congenital or acquired malformations in the cerebral or visceral circulation
\end{itemize}

Angiographic findings are always interpreted alongside clinical assessment and other imaging, such as CT, MRI, ultrasound, and functional tests of blood flow.

\section*{When to Seek Medical Care}
Contact the hospital or your doctor promptly after an angiogram if you develop:

\begin{itemize}
    \item Heavy bleeding, an enlarging bruise, or a painful, firm swelling at the puncture site
    \item Coldness, numbness, tingling, or a pale or blue colour in the limb used for the procedure
    \item Chest pain or breathlessness that does not settle
    \item Fever, chills, or spreading redness around the puncture site
    \item Difficulty passing urine, since the contrast dye can stress the kidneys
\end{itemize}

\textbf{Call 000 (or local emergency services) for:}
\begin{itemize}
    \item Chest pain or tightness suggesting a heart problem
    \item Sudden weakness, numbness, or difficulty speaking, which may indicate a stroke
    \item Heavy, uncontrollable bleeding from the puncture site
    \item Fainting, collapse, or sudden loss of consciousness
\end{itemize}

\section*{Diagnosis and Investigations}
Preparation for an angiogram includes a careful pre-procedure assessment:

\begin{itemize}
    \item \textbf{History and examination:} documentation of allergies, current medicines (especially blood thinners), kidney function, pregnancy status, and any previous contrast reactions
    \item \textbf{Blood tests:} kidney function (creatinine), clotting screen, and full blood count
    \item \textbf{ECG:} routine before coronary angiography
    \item \textbf{Informed consent:} a discussion of the benefits, risks, and alternatives, including non-invasive options
    \item \textbf{Fasting:} generally required for several hours beforehand, particularly if sedation is planned
\end{itemize}

During the procedure, X-ray screening guides the catheter under local anaesthetic. Afterward, the images are reviewed to identify narrowing, blockages, aneurysms, or bleeding. Additional measurements are sometimes made inside the vessel, including pressure gradients across a narrowing (fractional flow reserve) or ultrasound imaging from within the artery (intravascular ultrasound), to judge whether a stenosis is significant enough to require treatment.

\section*{Management}
An angiogram is part of diagnosis, and management depends on what is found:

\begin{itemize}
    \item \textbf{Normal results:} reassurance, with attention to cardiovascular risk factors where relevant
    \item \textbf{Mild narrowing:} lifestyle change and medical therapy rather than intervention
    \item \textbf{Significant narrowing:} a balloon and stent (angioplasty) can be placed in the same session
    \item \textbf{Aneurysm or bleeding:} the vessel may be sealed with coils, glue, or a stent graft during the same procedure
    \item \textbf{Aftercare:} lie still for several hours with the puncture limb kept straight, and use pressure or a closure device at the site
    \item \textbf{Hydration:} drink plenty of fluids to help flush the contrast dye through the kidneys
    \item \textbf{Follow-up:} review of the results with the referring team and onward referral to the relevant specialist
\end{itemize}

\section*{Complications}
\begin{itemize}
    \item Bruising and bleeding at the puncture site
    \item Arterial damage, including a false aneurysm (pseudoaneurysm) or a tear of the vessel wall (dissection)
    \item Contrast dye reactions, ranging from a mild rash to a severe anaphylactic response
    \item Contrast-induced kidney injury, especially with pre-existing renal disease or dehydration
    \item Blood clot formation, which can block a vessel and cause a heart attack, stroke, or limb ischaemia
    \item Infection at the puncture site
    \item Radiation exposure from the X-rays, kept as low as reasonably achievable
\end{itemize}

Serious complications are uncommon in experienced centres, but they are discussed with the patient before consent.

\section*{Prognosis and Outcomes}
For most people an angiogram is a safe procedure with a low risk of serious complications and rapid recovery; most return to normal activities within a day or two. The principal benefit is accurate anatomical information that guides treatment, and because diagnosis and treatment are frequently combined, many people leave the procedure with their condition treated rather than merely diagnosed. The long-term outlook is determined by the condition that was found and its treatment, not by the test itself, and it is generally excellent when the abnormality is managed appropriately.

\section*{Prevention and Self-Care}
\begin{itemize}
    \item Tell the team about any allergies, especially a previous reaction to contrast dye
    \item Inform them of all medicines, particularly blood thinners and diabetes medications
    \item Follow fasting instructions and arrange for someone to drive you home afterwards
    \item Keep the puncture site clean and dry for the first day or two
    \item Report any swelling, bleeding, numbness, or unusual pain promptly
    \item Drink plenty of fluids after the procedure to protect the kidneys
    \item Follow the advice given about activity, lifting, and returning to work
\end{itemize}

\section*{Sources and Further Reading}
The following sources provide reliable further information on angiography:

\begin{itemize}
    \item NHS. \textit{Information on coronary angiography and what it involves.} https://www.nhs.uk/conditions/coronary-angiography/
    \item MedlinePlus. \textit{Overview of angiography, including preparation and risks.} https://medlineplus.gov/angiogram.html
    \item British Heart Foundation. \textit{Plain-language guide to coronary angiography and angioplasty.} https://www.bhf.org.uk/informationsupport
    \item Mayo Clinic. \textit{Coronary angiogram: purpose, procedure, and risks.} https://www.mayoclinic.org/tests-procedures/coronary-angiogram/about/pac-20384904
    \item RadiologyInfo. \textit{Angiography and vascular imaging for patients, from the Radiological Society of North America.} https://www.radiologyinfo.org/
\end{itemize}
